\section*{Open Questions (Post-Replication)}

The following five questions remained open after this replication and are the
most directly actionable follow-ups. Each identifies a specific gap that our
replication did not close (either because the paper itself does not address it,
or because our scoped-in tests did not exercise it).

\begin{enumerate}

\item \textbf{Charge-transfer and Rydberg states.}
Does MC-VQE retain sub-100-$\mu$eV accuracy when extended to CT and Rydberg
excited states, which are notoriously poorly described by the CIS-based
reference-state construction MC-VQE relies on? \\
\emph{Basis:} MC-VQE builds contracted references from CIS eigenvectors. CIS is
known to badly under-estimate CT/Rydberg energies (often by 0.5--1 eV). If the
contracted subspace lacks CT character, the state-averaged entangler must
compensate --- but the paper (and our replication) only exercises
locally-excited exciton states where CIS is already qualitatively correct. \\
\emph{Next step:} Small donor-acceptor dimer (4--8 qubits, STO-3G) with a known
CT root; compare MC-VQE-CIS vs.\ MC-VQE-CIS(D) vs.\ FCI.

\item \textbf{Head-to-head vs.\ modern excited-state VQE variants.}
How does MC-VQE compare to qEOM (McClean/Kim 2020), ADAPT-VQE-excited, and
subspace-search VQE on identical molecules, shot budgets, and noise models? \\
\emph{Basis:} The paper predates qEOM and ADAPT-based excited-state schemes. No
published head-to-head resource comparison exists on a common benchmark. \\
\emph{Next step:} H2, LiH, H4-chain at STO-3G, all four methods, common shot
budget, chemical accuracy for lowest 3 states. Report circuit depth, shots,
and depolarizing-noise robustness at $p=10^{-3}$--$10^{-2}$.

\item \textbf{Noise robustness of the subspace diagonalization step.}
What is MC-VQE's noise robustness on real hardware or realistic noise
simulation, and does classical diagonalization of the measured
$H_{\Theta\Theta'}$ matrix amplify or suppress shot noise? \\
\emph{Basis:} Shot noise on individual matrix elements propagates non-trivially
through the eigendecomposition; near-degenerate eigenstates can rotate strongly.
The paper is noiseless, as is our replication. \\
\emph{Next step:} Add shot-based measurement layer to \texttt{mcvqe\_exciton.py};
sweep $10^3$--$10^6$ shots/element under Qiskit-Aer depolarizing +
thermal-relaxation noise. Metric: shot budget for chemical-accuracy excitation
energies.

\item \textbf{Conical intersections and photochemistry.}
Can MC-VQE be applied at conical intersections, and does its simultaneous
multi-state formulation cleanly resolve state-crossings where sequential VQD
becomes unstable? \\
\emph{Basis:} MC-VQE's philosophy is well-matched to CoIns, but the paper only
tests the isolated LH2 B850 ring far from crossings. This is the most
scientifically important untested regime, since photochemistry is where
excited-state QC should matter most. \\
\emph{Next step:} 1D geometry scan across a known CoIn (H3, pyrrole N-H,
formaldehyde) at STO-3G/6-31G; compare MC-VQE energy curves to MRCI/CASSCF for
continuity, absence of spurious avoided crossings, correct branching-space
recovery.

\item \textbf{Scaling to real chromophore active spaces.}
Does MC-VQE's ``single entangler layer suffices'' scaling survive when moving
from the structured exciton Hamiltonian (dipole-dipole NN couplings) to a
general CAS Hamiltonian for a real chromophore (retinal PSB, GFP,
pigment-protein complex)? \\
\emph{Basis:} Our $N=4$ tests already showed the ``single layer'' claim is
system-specific ($L=1$ failed by 748 meV; $L=3$ was needed for $25.6\,\mu$eV).
Real CAS Hamiltonians have dense two-electron integrals unlike the exciton
model's sparse Pauli structure. \\
\emph{Next step:} Ethylene CAS(2,2)$\to$(4,4)$\to$(6,6) then retinal PSB
CAS(6,6), track $L(\text{qubits})$ for chemical accuracy. If the paper's
$L=1$ scaling breaks, that is a substantive finding.

\end{enumerate}
